\section*{Chapitre \ref{ch:g5:people-and-nature} --- Les humains et la nature}

\begin{solution}{exo:g5:people-and-nature:1}
Les champs, gérés pour une culture choisie~; les \omterm{prop:g5:people-and-nature:managed}{forêts}, gérées pour
le bois~; les prés, gardés ouverts pour le foin ou le pâturage~; les
villes, les routes et les barrages, bâtis pour habiter et se
déplacer. (Quatre au choix, avec leur but.)
\end{solution}

\begin{solution}{exo:g5:people-and-nature:2}
Laissé sans fauche ni pâture, il se referme en broussailles en
quelques années, puis en bois --- le milieu ouvert et fleuri du pré
disparaît.
\end{solution}

\begin{solution}{exo:g5:people-and-nature:3}
Ne pas prélever plus vite que cela ne se renouvelle~; garder les
ouvriers (les \omterm{def:g5:biodiversity:species}{espèces} dont les services gratuits font vivre
l'usage)~; rendre ce qu'on emprunte (refermer les boucles, ne
laisser aucun déchet s'entasser).
\end{solution}

\begin{solution}{exo:g5:people-and-nature:4}
Haies arrachées ou conservées~: A perd ceux qui retiennent la terre
et les maisons des mangeurs de ravageurs. Terre nue ou couverte en
hiver~: celle de A part vers le bas avec les pluies. Une seule
culture géante ou une rotation~: A invite n'importe quel ravageur de
cette culture à balayer tout le coteau, et ne repose jamais le sol.
\end{solution}

\begin{solution}{exo:g5:people-and-nature:5}
Compter ce que la forêt produit en une année~; n'en couper pas
davantage. C'est la règle 1 --- ne pas prélever plus vite que cela
ne se renouvelle --- et la forêt donne du bois chaque année,
indéfiniment, tout en restant une forêt vivante.
\end{solution}

\begin{solution}{exo:g5:people-and-nature:6}
Par exemple~: les martinets, qui nichent sur des «~falaises~» de
béton~; les renards, qui patrouillent poubelles et jardins~; les
\omterm{def:g1:plants-around-us:plant}{plantes} sauvages, enracinées dans les joints du bitume et des murs
(aussi les pigeons sur les corniches, les faucons sur les tours).
\end{solution}

\begin{solution}{exo:g5:people-and-nature:7}
Pour les humains~: une eau ralentie et des crues stockées dans les
prairies humides restaurées --- moins d'inondations en aval. Pour le
réseau~: les milieux des méandres sont revenus, et avec eux les
\omterm{prop:g3:sorting-living-things:groups}{poissons}, les libellules et les prés de chasse des cigognes.
\end{solution}

\begin{solution}{exo:g5:people-and-nature:8}
Réponse libre. Trouvailles habituelles~: bas-côtés fauchés, arbres
taillés, massifs plantés, trottoirs désherbés --- avec peut-être un
haut de mur ou un fossé non géré, souvent l'endroit le plus sauvage
de la vue.
\end{solution}

\begin{solution}{exo:g5:people-and-nature:9}
Ne jamais y toucher détruirait en fait certains milieux~: le pré
fleuri existe \emph{parce qu'}on le fauche, et se referme sans cela.
Protéger la nature n'est pas ne pas y toucher, mais bien la gérer
--- la question honnête est comment toucher~: quelles règles l'usage
respecte, et non si des mains s'y sont posées.
\end{solution}

\begin{solution}{exo:g5:people-and-nature:10}
La règle 2 franchement enfreinte (les ouvriers expulsés avec les
haies), la règle 3 en partie (une terre emportée n'est jamais
rendue), et la règle 1 menacée (la terre se perd plus vite qu'elle
ne se forme). Le changement unique le plus efficace~: replanter les
haies --- elles retiennent la terre et relogent les mangeurs de
ravageurs d'un seul coup.
\end{solution}

\begin{solution}{exo:g5:people-and-nature:11}
Les prises doublées ont pris les \omterm{prop:g3:sorting-living-things:groups}{poissons} bien plus vite que les
bancs ne se reproduisaient --- en entamant le stock reproducteur
lui-même. Dépenser son capital paraît toujours magnifique tant que
cela dure~: les cales étaient pleines précisément parce que les
\omterm{prop:g3:sorting-living-things:groups}{poissons} de l'an prochain étaient débarqués cette année. Le
forestier ne récolte que la pousse de l'année et garde le capital
sur pied intact --- même arithmétique, signe opposé, et c'est
pourquoi la forêt donne pour toujours et la flotte s'est arrêtée la
sixième année.
\end{solution}
